% Source: https://artofproblemsolving.com/community/c3h1423932_easter_math_competition_series
% Stable identity: aops-topic-1423932-q07
Let $\alpha$ be half of the angle at $A$. Both circles are tangent to the two sides from $A$, so their centres lie on the bisector of angle $A$ (the altitude from $A$), and a circle of radius $\rho$ tangent to both sides has its centre at distance $\dfrac{\rho}{\sin\alpha}$ from $A$. The circles are externally tangent, with $P$ closer to $A$, so
\[
\dfrac{x}{\sin\alpha}-\dfrac{1}{\sin\alpha}=x+1\qquad\Longrightarrow\qquad \sin\alpha=\dfrac{x-1}{x+1},\qquad \cos\alpha=\dfrac{2\sqrt x}{x+1}.
\]
The triangle has half-base $6\sqrt2\sin\alpha$ and height $6\sqrt2\cos\alpha$, so its inradius (area divided by semiperimeter) is
\[
x=\dfrac{\left(6\sqrt2\right)^2\sin\alpha\cos\alpha}{6\sqrt2+6\sqrt2\sin\alpha}=\dfrac{6\sqrt2\sin\alpha\cos\alpha}{1+\sin\alpha}.
\]
With $1+\sin\alpha=\dfrac{2x}{x+1}$ this becomes
\[
x=6\sqrt2\cdot\dfrac{x-1}{x+1}\cdot\dfrac{2\sqrt x}{x+1}\cdot\dfrac{x+1}{2x}=\dfrac{6\sqrt2\,(x-1)}{\sqrt x\,(x+1)},
\]
that is, $x^{\dfrac{3}{2}}(x+1)=6\sqrt2\,(x-1)$, or after squaring, $x^3(x+1)^2=72(x-1)^2$. We need $x>1$ (as $\sin\alpha>0$). Testing integers: $x=2$ gives $8\cdot9=72=72\cdot1$. For $x=3$ and $x=4$ the two sides are $432\ne288$ and $1600\ne648$; for $x\ge5$, $x^3(x+1)^2>125(x-1)^2>72(x-1)^2$. So $x=\boxed{2}$.

Check: $\sin\alpha=\dfrac{1}{3}$, so the triangle has half-base $2\sqrt2$, height $8$ and inradius $\dfrac{2\sqrt2\cdot8}{6\sqrt2+2\sqrt2}=2$; the centres are at heights $2$ and $5$ above $BC$, $3$ apart, and the upper circle is at distance $3\cdot\dfrac{1}{3}=1$ from the sides.

\begin{asy}
usepackage("amsmath");
size(5.5cm);
defaultpen(fontsize(10pt));
pair A=(0,8),B=(-2*sqrt(2),0),C=(2*sqrt(2),0),O=(0,2),P=(0,5);
draw(A--B--C--cycle,linewidth(0.9));
draw(A--(0,0),gray+dashed);
draw(circle(O,2),blue+linewidth(0.9)); draw(circle(P,1),red+linewidth(0.9));
pair TO=B+dot(O-B,unit(A-B))*unit(A-B);
pair TP=B+dot(P-B,unit(A-B))*unit(A-B);
draw(O--TO,blue+dashed); draw(P--TP,red+dashed);
dot(A);dot(B);dot(C);dot(O,blue);dot(P,red);
label("$A$",A,N); label("$B$",B,SW); label("$C$",C,SE); label("$O$",O+(0.45,0.3),blue+fontsize(9pt)); label("$P$",P+(0.38,0.22),red+fontsize(9pt));
label("$x$",(O+TO)/2+0.4*unit(rotate(90)*(O-TO)),blue+fontsize(9pt)); pair L1=(-1.6,6.1); draw(L1--(P+TP)/2,red+linewidth(0.3)); label("$1$",L1,W,red+fontsize(9pt));
label("$6\sqrt{2}$",(A+C)/2,NE);
label("\mbox{}",(min(currentpicture,user=true)+max(currentpicture,user=true))/2,black); // keeps the text after the figure black when LaTeX typesets the labels inline
\end{asy}
